\chapter{La capa fina}

\label{cap:opacidad}
\section{Un obstáculo que se repite}

Este capítulo no estaba previsto. Se impuso porque el mismo obstáculo apareció cinco veces, en instrumentos de dos niveles de gobierno distintos ---el municipal y el provincial---, siempre en el mismo lugar del procedimiento: el punto donde el acto pasa de enunciar a decidir.

\textbf{Una corrección que el propio capítulo debe poner al frente.} Cuatro de esos cinco anexos resultaron estar, después, en el repositorio del Boletín: uno escaneado con capa de texto y tres como imágenes escaneadas sin ella. El quinto, el reglamento de fraccionamiento de Vaqueros, apareció al final fuera del Boletín: en el digesto que el propio municipio publica en línea. Lo que este capítulo documenta son, entonces, dos cosas que al principio parecían una: \textbf{anexos que no se publican} y \textbf{anexos que se publican de un modo que nadie los encuentra}.

Lo que sigue es el inventario --- seis instrumentos, de dos niveles de gobierno, relevados a lo largo de todo este trabajo --- y después la interpretación.

\section{El departamento se catastró en 1929}

Este capítulo sostiene que la opacidad de este territorio está en el criterio y no en el
acto. Conviene empezar por establecer desde cuándo existe el registro que hoy no se puede
consultar sin arancel.

\begin{ficha}{Decreto Nº 10.910 del 5 de agosto de 1929}
Expediente 3463-D. Vista la nota del Director de la Comisión de Catastro elevando la
solicitud del señor \textbf{Adeodato Aybar\index{Aybar, Adeodato}}, «en la cual éste manifiesta que teniendo
urgente necesidad de la suma de \$500, y \textbf{habiendo terminado la catastración de los
Departamentos de Campo Santo, Caldera y La Viña}, pide se le liquide la citada suma a
cuenta de la remuneración que le sea fijada».

Se liquidan quinientos pesos al Director de la Comisión de Catastro para que los entregue
al catastrador, imputados a la Ley 9720.
\end{ficha}

De modo que el departamento La Caldera está catastrado desde 1929, y el trabajo lo hizo una
persona con nombre y apellido, a la que el Estado le debía la remuneración cuando terminó.

Ese dato tiene dos consecuencias para este libro.

La primera es que \textbf{el registro existe desde hace casi un siglo}. Cuando el capítulo \ref{cap:fincas} muestra mensuras que describen límites por «unos mojones antiguos», y cuando este libro
no puede establecer si una matrícula proviene de otra, no está describiendo un territorio
sin catastrar: está describiendo un territorio catastrado en 1929 cuyo catastro no es
consultable sin arancel ni presencia física.

La segunda es de escala, y explica algo del carácter de la información que quedó. Un solo
catastrador terminó tres departamentos y cobró quinientos pesos a cuenta. El presupuesto
municipal de La Caldera de 1908 era de mil cuatrocientos ochenta y tres; la casa quinta de
Lesser se remató en 1927 con base de seis mil. El levantamiento catastral de tres
departamentos costó, en pagos a cuenta, menos de la décima parte de una finca de diez
hectáreas.

\section{Los anexos sí se publican}

Este capítulo sostiene que seis instrumentos publican la envoltura y reservan el contenido
decisorio a un anexo. Corresponde precisar esa afirmación antes de desarrollarla, porque el
Boletín Oficial tiene, desde hace años, un repositorio de anexos.

\textbf{Existe una colección digital de anexos, separada de la del boletín.} Un relevamiento
de dieciséis anexos de octubre de 2017 arroja \textbf{trescientas veintiséis páginas} de
material: transferencias de partidas presupuestarias, incorporaciones de obra pública,
planillas de personal. La Caldera y Vaqueros aparecen ahí --- con código auxiliar, importe y
número de expediente --- en lugares donde el acto publicado en el boletín no los nombra.

De modo que \textbf{la afirmación de este capítulo no puede ser que los anexos no se
publiquen}. Se publican, y en volumen.

Lo que este libro documenta es más específico, y conviene enunciarlo con exactitud:

\begin{quote}
De seis instrumentos identificados cuyo contenido operativo está en un anexo o en la sede del organismo, \textbf{cinco no se obtuvieron por ninguna vía}, incluida la consulta
del repositorio. El sexto, el reglamento de fraccionamiento de Vaqueros, lo terminó publicando el propio municipio en su digesto en línea, no el Boletín que publicó el acto.
\end{quote}

La diferencia importa. No es que el Estado no publique anexos: es que \textbf{unos anexos se
publican y otros no, y no hay criterio visible que distinga a unos de otros}. Un anexo de
transferencia de partidas por quince mil quinientos noventa y cinco pesos con siete centavos
está disponible; el contrato de promoción turística que fija exenciones impositivas, no.

\begin{cautela}
Hay además un problema de acceso que conviene separar del de publicación, porque son
distintos y este libro los ha tratado juntos.

En ese mismo relevamiento de dieciséis archivos, \textbf{ninguno es un PDF pese a llevar esa
extensión}: trece son archivos comprimidos con una imagen y un texto por página, dos son
texto plano corrido sin marcas de página --- lo que impide citar por página cualquier
hallazgo que contengan --- y uno pesa trescientos catorce bytes, todos saltos de línea.
\textbf{Dos de los trece venían con la capa de texto enteramente vacía} y hubo que
reconocerlos ópticamente para poder leerlos.

Nada de eso es ocultamiento. Es infraestructura documental deficiente, y produce el mismo
efecto práctico: \textbf{el contenido está publicado y no es consultable sin trabajo técnico
previo}. Este libro lo consigna como lo que es --- una barrera de acceso, no una decisión de
reserva --- y lo distingue de los cinco contenidos que no se obtuvieron por ninguna vía.
\end{cautela}

\subsection*{Y un caso distinto: el acto y el acta, sin el detalle}

Los seis instrumentos de este capítulo publican el acto y reservan el anexo. El Decreto
1.889/16 publica las dos cosas, y conviene anotarlo porque completa el cuadro.

La ficha web del Boletín (Nº 19.913, 29/11/2016) trae el decreto: visto, considerandos,
articulado y refrendo. El repositorio de anexos ofrece, con el nombre del decreto, el
\textbf{Acta Acuerdo} del Plan de Mínimas y Encauzamiento 2016--2017 --- cláusulas,
mecanismo de pago, firmantes ---, escaneada desde el folio 3. Ninguna de las dos piezas
transcribe la fecha del decreto ni el número de expediente.

\textbf{De modo que se conocen el acto y el convenio, y no su contenido operativo.} Dentro
del convenio, las dos piezas que deciden --- los proyectos de obra y la distribución de fondos
por municipio --- se invocan como existentes y no se adjuntan.

Es la misma estructura en otro orden: \textbf{lo publicado permite saber que algo se decidió, y
no qué se decidió ni sobre qué río.}

\section{Seis instrumentos con el criterio afuera}

\begin{ficha}{Instrumentos cuyo contenido operativo no se publicó junto al acto}
\begin{itemize}[leftmargin=*,nosep]
\item \textbf{Ordenanza Nº 000437, Municipalidad de Vaqueros} (Boletín Oficial provincial, 15/01/2015). Aprueba el Reglamento de Fraccionamiento de Tierras. Publicado: cuatro artículos aprobatorios. No publicado: el reglamento, que contiene superficies mínimas, cesiones obligatorias, obras exigibles al loteador, procedimiento y condiciones ambientales. Al pie, el Boletín consigna que el anexo se encuentra para su consulta en sus propias oficinas («de esta reparación», dice el original), y la ficha web del instrumento no ofrece enlace «Ver anexo». Es el único de los cinco casos con que este capítulo empezó cuyo anexo no está en el repositorio del Boletín. \textbf{Obtenido en septiembre de 2026 del Digesto Jurídico Virtual de la Municipalidad de Vaqueros}, una carpeta pública en línea que el municipio ya anunciaba el 2 de febrero de 2024; su contenido se resume en el capítulo \ref{cap:vaqueros}.
\item \textbf{Resolución Nº 383/14, Secretaría de Recursos Hídricos} (edicto en los B.O.\ Nº 19.466 y 19.467, 21 y 22/01/2015; la ficha web no ofrece enlace «Ver anexo»). Aprueba la determinación de la línea de ribera en ambas márgenes de los arroyos El Durazno y Guaranguay, matrícula 4366, departamento La Caldera. Publicado: que la línea fue determinada. No publicado: \textbf{las coordenadas Posgar que indican dónde está}, que quedan ``a disposición de los interesados en la sede de la Secretaría''. Otro expediente de la misma sección, en la misma edición del Boletín, sí publica su tabla de coordenadas. \textbf{No obtenidas.}
\item \textbf{Decreto Nº 463/15} (09/02/2015; B.O.\ Nº 19.485). Ratifica el contrato de promoción turística con Proyecto 1 S.A.\ para ampliar el Hotel La Caldera. Publicado: el acto ratificatorio. \textbf{No publicado: el contrato}; al pie, la leyenda «anexo para su consulta en oficinas de nuestra repartición», la misma fórmula de la Ordenanza 437. No obtenido.
\item \textbf{Decreto Nº 3069/14} (14/10/2014). Ratifica el contrato de promoción turística con La Caldera S.A. Publicado: el acto ratificatorio. \textbf{No publicado: el contrato}, que fija los alcances de las exenciones impositivas concedidas; la ficha del decreto en la edición web del Boletín no ofrece enlace «Ver anexo». No obtenido.
\item \textbf{Resolución Nº 142/14}, Secretaría de Recursos Hídricos. Aprueba la determinación de línea de ribera de los ríos Mojotoro y La Caldera sobre la matrícula 3616, y la Resolución 236/14 la rectifica. Publicado: el acto y el inmueble. \textbf{No publicadas: las coordenadas}, ni las erróneas ni las corregidas, remitidas a la sede del organismo (avisos en los B.O.\ Nº 19.364 y 19.377, agosto y septiembre de 2014, sin enlace «Ver anexo»). Con la 383/14, son las dos únicas de las trece determinaciones de 2013--2015 referidas a inmuebles particulares ---las quince de la serie, menos la 467/15, que delimita un tramo de curso, y la 023/14, que recae sobre el pueblo y no sobre una matrícula--- que no las publicaron, y ambas recaen sobre suelo en proceso de loteo, aunque esa condición la comparten otras dos que sí las publicaron (capítulo \ref{cap:ribera}). No obtenidas.
\item \textbf{Informe técnico del inciso e}, que funda la contratación por vía abreviada del Plan de Manejo. Publicado: el encuadre legal. \textbf{No publicado: el informe}. No obtenido.
\end{itemize}
\end{ficha}

\begin{cautela}
De los seis, cinco no se consiguieron. Los que este libro sí consiguió resultaron estar publicados, y el sexto también: no donde lo publicó el acto, sino en el digesto del municipio que lo dictó. Desde 2018 la Ley 8126 obliga a los municipios a publicar sus ordenanzas en el Boletín Oficial y en internet (art.\ 79); Vaqueros lo cumple por su propia cuenta, el Boletín provincial sigue remitiendo a sus oficinas.

\textbf{La lista tenía trece, y ocho salieron de ella.} Los que quedan fueron verificados uno por uno en la edición web: ninguno ofrece enlace «Ver anexo». La nómina de adjudicatarios del Decreto 63/25, el Reglamento de Asignación de la Resolución 82/16, el baremo de la Resolución 8/16 y los procedimientos de la Resolución 9/16 fueron los primeros casos que este libro registró como anexos no publicados, y los obtuvo por otras vías. Estaban publicados: la edición web del Boletín ofrece cada uno como archivo vinculado al acto, bajo la leyenda «Ver anexo». Los tres del régimen de tierra fiscal, además, son copias certificadas escaneadas, sin capa de texto, que ningún buscador encuentra; con los de la Resolución 7/16, que no estaba en la lista, y los de la 567D/26, los anexos de ese régimen publicados sólo como imagen son cinco. El del Decreto 63/25 es un caso intermedio: también es una copia certificada escaneada, de cuatro folios, pero trae una capa de texto de reconocimiento óptico ---imperfecta, con números de expediente y documentos mal leídos--- que lo vuelve buscable. A ellos se sumaron después el anexo del Decreto 1682/19, que reglamenta la aprobación de todos los loteos de la provincia ---también vinculado al acto, y con texto legible---, el de la Resolución 023/14 de Recursos Hídricos, que resultó contener la línea de ribera del pueblo de La Caldera con sus 134 coordenadas, en una copia escaneada, los de la Resolución 567D/26, que este libro contaba como obtenidos por vía lateral y estaban vinculados al acto, y el convenio de licencias de conducir del Decreto 1433/15, también escaneado. Los ocho pasan a la columna de lo que el repositorio sí contiene ---los escaneados, a la de la barrera de acceso que el apartado anterior describe---, y confirman su advertencia: la frontera entre lo publicado y lo reservado no siempre está donde un lector la busca, y cada uno de los cinco restantes debe cotejarse contra ese repositorio antes de citarse como no publicado. El cotejo sumó, en cambio, un caso: el Decreto 463/15, gemelo del 3069/14, cuya ficha remite el contrato a las oficinas del Boletín con la misma fórmula que la Ordenanza 437. Con él, la lista queda en seis. Queda fuera, porque no es un anexo sino un acto entero, la Resolución SMyE 20/2022 que rige la zona de exclusión de canteras junto a las tomas de agua: no se publicó (capítulo \ref{cap:amparo}).

Los anexos no son secretos: existen, integran los actos y son accesibles por pedido o por vías laterales. Lo que ocurre es que la publicidad legal se cumple sobre la envoltura y la publicidad efectiva queda a cargo del interesado. Quien tiene tiempo, abogado o capacidad de gestión los obtiene. Quien lee el Boletín Oficial de buena fe, no.

Y el contenido reservado es siempre el mismo tipo de contenido: qué obras debe hacer un loteador, cuántos puntos vale una discapacidad frente a una familia numerosa, cuántos días tiene un vecino para impugnar un padrón, quién recibió cada parcela. Es decir, aquello que permitiría a un ciudadano saber si fue tratado con la misma vara que otro.
\end{cautela}

\begin{cautela}
El caso de la Resolución 82/16, que este libro presentaba como el ejemplo central del patrón, muestra en cambio la otra forma del problema: la barrera de acceso. Es la cabeza de la cadena y todas las demás remiten a ella, y su anexo está publicado, pero sólo como imagen.

La Resolución 9/16 remite tres veces el plazo de exhibición de padrones a la Resolución 82/16, dos de ellas a su artículo 5º (capítulo \ref{cap:tierrafiscal}).

Ese plazo es el número del que depende toda la arquitectura de control ciudadano del régimen de tierra fiscal, y no está en ninguna de las resoluciones que se publican como texto. El anexo, una copia certificada de seis folios escaneada sin capa de texto, no aparece en ninguna búsqueda: hay que abrirlo y leerlo. Su artículo 5º no habla de padrones: regula la publicidad de los sorteos y fija un solo plazo, diez días de anticipación. Por remisión, entonces, los padrones deben exhibirse al menos diez días. Se publican las garantías como texto, y el plazo que las hace ejercitables queda en una imagen y hay que deducirlo de una cláusula escrita para otra cosa.
\end{cautela}

\section{La excepción que se tramita}

El segundo mecanismo es distinto y opera en sentido complementario. No consiste en no publicar el criterio, sino en publicar una prohibición que la propia norma convierte en trámite.

El Código de Ordenamiento Territorial de La Caldera declara zonas no urbanizables --- entre ellas la Zona PU, sobre ambas márgenes del río --- y prevé al mismo tiempo la figura del Proyecto Especial del artículo 31, que somete el caso a dictamen del Órgano Técnico de Aplicación y aprobación por ordenanza. En diciembre de 2020, ante el loteo La Falda, el municipio determinó que el sector no es urbanizable porque la urbanización degradaría el ambiente, y en el mismo acto indicó al interesado que presentara un proyecto especial ajustado al Código con certificado provincial, para que el Concejo lo analizara y aprobara.

El proyecto de ordenanza del PDUA deja seis decisiones «a consideración del municipio» sin decir qué debe evaluarse, con qué dictamen, en qué plazo ni con qué recurso (capítulo \ref{cap:pdua}).

La calificación de no urbanizable no opera como prohibición sino como requisito de procedimiento agravado. Lo que en el plano de zonificación es un límite, en el trámite es una puerta con más pasos.

Nada de esto es irregular. El capítulo \ref{cap:cot} ya mostró la paradoja del Proyecto Especial: la figura existe para desarrollos que requieren tratamiento particularizado, y el municipio carece de aparato para particularizar nada. Aquí interesa sólo su efecto sobre lo que se puede saber.

\section{Un mapa que no dice dónde}

El caso más nítido de las tres capas apareció al final del relevamiento, es provincial, y es de 2025.

\begin{ficha}{Ley provincial Nº 8483 y su reglamentación}
\begin{itemize}[leftmargin=*,nosep]
\item \textbf{Ley Nº 8483}, publicada en el Boletín Oficial provincial el 3 de enero de 2025: revisa el Ordenamiento Territorial de Bosques Nativos y modifica las superficies por categoría de conservación. Es la norma que el senador por La Caldera votó en contra en diciembre de 2024.
\item Introduce el \textbf{Área de Producción y Conservación (APC)}, categoría transitoria, y el concepto de \textbf{área verde}.
\item Modifica el artículo 14 de la Ley 7543 y habilita en Categoría II los Sistemas de Ganadería Silvo-Pastoril y el Manejo de Bosque con Ganadería Integrada.
\item Prevé que las áreas comprendidas en las categorías de conservación puedan ser objeto de certificación voluntaria de \textbf{bonos verdes} vinculados al cambio climático, siempre que cuenten con plan de conservación o de manejo aprobado.
\item \textbf{Resolución SAyDS Nº 00509}, B.O.\ 27/08/2025: los inmuebles rurales con sectores de APC podrán ser objeto de autorización de cambio de uso de suelo, identificando dentro del APC la Categoría III sobre la cual se establecerán los proyectos, y quedando el remanente en Categoría II.
\item \textbf{Disposición nacional Nº 2483/2025}, B.O.\ 20/10/2025: acredita la actualización. Superficies resultantes: Categoría I (rojo) 1.308.244 ha; Categoría II (amarillo) 5.528.753 ha; Categoría III (verde) 721.568 ha; total de bosque nativo 7.558.565 ha.
\item \textbf{Resolución SAyDS Nº 00235}, B.O.\ 13/04/2026: aprueba el reglamento de requisitos para proyectos de intervención sobre bosques nativos y deroga las Resoluciones 558/13, 831/19, 333/20 y 411/20.
\end{itemize}
\end{ficha}

\subsection*{Cómo se hizo el mapa}

El Anexo I de la Ley 8483 contiene el informe final del equipo técnico, y su lectura obliga a corregir la impresión que la sola descripción del mecanismo produce.

\begin{ficha}{Proceso de revisión del OTBN de Salta --- Anexo I de la Ley 8483}
\begin{itemize}[leftmargin=*,nosep]
\item \textbf{Equipo técnico:} convocado por Decreto provincial 3749/2014 e integrado por la Secretaría de Ambiente y Desarrollo Sustentable, el INTA, el INENCO-CONICET, la Facultad de Ciencias Naturales de la Universidad Nacional de Salta, el INAI y la Administración de Parques Nacionales. Informe final del 30 de noviembre de 2023.
\item \textbf{Marco:} artículo 6 de la Ley nacional 26.331 y Decreto 91/2009, que obligan a actualizar el ordenamiento cada cinco años. Se declara respetado el principio de progresividad del artículo 4 de la Ley General del Ambiente.
\item \textbf{Herramientas:} dos programas de acceso libre y gratuito desarrollados en Google Earth Engine. El \textbf{Visor}, repositorio donde puede consultarse y descargarse toda la información utilizada. El \textbf{Mapeador}, que elabora el mapa mediante evaluación multiobjetivo en tres pasos.
\item \textbf{Taller participativo:} 25 de octubre de 2023, con un Consejo Asesor de 53 instituciones y 91 participantes acreditados, agrupados en cuatro perspectivas identificadas por metodología Q: respeto a la diversidad cultural; equilibrio entre Estado y mercado; producción agropecuaria sin deforestación; y el campo como motor de la economía.
\item \textbf{Decisión declarada como inédita:} la participación pública fue \textbf{vinculante}. Los mapas surgidos del taller fueron el único insumo del mapa final. El propio informe consigna que las leyes vigentes no obligaban al Estado a hacerlo.
\item \textbf{Superficies del informe técnico:} Categoría I 1.278.221 ha (16,95\%); Categoría II 5.539.750 ha (73,48\%), de las cuales 3.247.625 en amarillo estricto y 2.292.125 en amarillo potencial; Categoría III 721.568 ha (9,57\%); APC 3.013.692 ha (39,97\%). Las áreas verdes equivalen al 23,94\% del APC. \textit{Las superficies de las Categorías I y II difieren de las que acredita la Disposición nacional 2483/2025 (1.308.244 y 5.528.753 ha); la diferencia la explica la adenda del Anexo I: la autoridad provincial ajustó la cartografía del informe ---filtró polígonos pequeños e incorporó bosques del Monte--- para responder a las observaciones de la autoridad nacional, y las superficies acreditadas son las de esa adenda.}
\end{itemize}
\end{ficha}

\begin{cautela}
Corresponde una corrección expresa de este capítulo.

La sola descripción del mecanismo --- una categoría que autoriza el desmonte y carece de ubicación geográfica --- sugiere ocultamiento. El informe técnico muestra lo contrario: un proceso con herramientas públicas y descargables, con participación vinculante de cincuenta y tres instituciones, con metodología declarada, con equipo técnico interinstitucional que incluye al CONICET, la universidad pública, Parques Nacionales y el INAI, y con una decisión política explícita de someterse a un resultado participativo que ninguna ley le exigía.

Y el APC tiene fundamentos técnicos enunciados. El informe sostiene que distribuir el verde dentro de un área transitoria, en lugar de fijarlo en polígonos, \textbf{permite minimizar los impactos puntuales y acumulativos} de los proyectos futuros y facilita el establecimiento de corredores ecológicos; y que permite una asignación equitativa de zonas verdes entre los productores, en lugar de que la categoría beneficie a quien tuvo la suerte de quedar dentro del polígono. Establece además que las limitaciones por cuenca se vinculan a su historia productiva: \textbf{a mayor deforestación pasada, mayores restricciones futuras}.

Ninguno de esos argumentos es trivial y este libro no está en condiciones de refutarlos.
\end{cautela}

\textbf{Y el mismo informe trae, para esta cuenca, una cifra que ordena todo lo demás.} Su Anexo I asigna a la cuenca Mojotoro--Lavayén ---que contiene la microcuenca del río La Caldera--- 29.744 hectáreas de Área de Producción y Conservación y \textbf{cero coma cero por ciento categorizable a verde}. El criterio del propio informe es explícito: las cuencas excedidas en su umbral máximo de deforestación, calculado por el método del Número de Curva, no pueden usarse para proyectos de cambio de uso del suelo, «y por lo tanto la deforestación estará prohibida en esas cuencas». La Tabla 17 del informe (página 89) sombrea las cuencas excedidas, y la de Mojotoro--Lavayén es una de las seis sombreadas ---con De las Chuñas, Juramento Inferior, Rosario--Horcones, Salí y Urueña, las mismas que tienen cero por ciento de verde---. La adenda que la autoridad provincial agregó al informe para responder a las observaciones de la autoridad nacional repite la cifra ---29.744 hectáreas de APC y cero por ciento categorizable a verde--- y aclara que el APC no es una categoría de conservación: mientras no se defina, se fiscaliza como si fuera Categoría II.

\textbf{Y la cuenca cubre tres cuartos del departamento.} Las capas que el informe técnico ofrece en su repositorio ---la de cuencas y el ráster del mapa final--- permiten el cruce que el Anexo II, a escala provincial, no deja ver. Sobre el límite que forman los veintiocho radios del Censo 2022, unas 107.600 hectáreas, la cuenca Mojotoro--Lavayén ocupa \textbf{84.500 hectáreas, el 79 por ciento}; la de Grande--Perico, 3.300; y 19.900 ---el 18 por ciento, casi todas en el noroeste--- quedan fuera de la capa de cuencas. Dentro del departamento, el mapa asigna unas 28.600 hectáreas a la Categoría I, 45.500 a la II y 3.100 al Área de Producción y Conservación; de estas últimas, \textbf{3.050 ---entre el 98 y el 99 por ciento--- están en la cuenca excedida}, donde el cambio de uso del suelo queda vedado. \textbf{En la práctica, todo el bosque nativo que el ordenamiento reconoce en el departamento está bajo un régimen que prohíbe el desmonte}: por su categoría o por su cuenca (lámina \ref{fig:otbn}).

\begin{figure}[p]
\centering
\includegraphics[width=\textwidth,height=0.84\textheight,keepaspectratio]{fig-otbn-cuenca.png}
\caption[El ordenamiento de bosques en el departamento La Caldera]{\textbf{El ordenamiento de bosques en el departamento La Caldera.} A: categorías del mapa final del informe técnico de la Ley 8483 sobre el límite que forman los radios del Censo 2022, con sus fracciones y radios, el límite de la cuenca Mojotoro--Lavayén y, rayado, el radio 0101, que el mapa deja sin categoría. B: el pueblo de La Caldera, con la línea de ribera que la Resolución 023/14 determinó sobre ambas márgenes del río, trazada con sus 134 coordenadas. En blanco, los sectores sin categoría. Elaboración propia con las capas del repositorio del informe técnico (descargadas el 17 de septiembre de 2026), la cartografía de radios del Censo 2022 (INDEC) y las coordenadas del anexo de la Resolución 023/14.}\label{fig:otbn}
\end{figure}

\begin{cautela}
Tres precisiones sobre este cálculo.

\textbf{La capa es la del informe técnico, y reproduce su Tabla 17.} Para la cuenca Mojotoro--Lavayén, el ráster da 56.309 hectáreas rojas, 180.338 amarillas y 29.571 de Área de Producción y Conservación, contra 56.037, 180.784 y 29.744 de la tabla: diferencias menores al uno por ciento, propias del cálculo por píxel. No incorpora los ajustes de la adenda provincial, que filtró polígonos pequeños y sumó bosques del Monte.

\textbf{El resultado no depende del límite.} La superficie censal no coincide con las otras que este libro registra para el departamento (capítulo \ref{cap:fincas}), y ninguna es un deslinde oficial. Tres contornos alternativos dan resultados parecidos: la capa vigente de departamentos del Instituto Geográfico Nacional ---unas 109.200 hectáreas, prácticamente la superficie que el INDEC usa para el censo--- deja el 78 por ciento en la cuenca; el del mismo organismo a escala 1:250.000 ---unas 110.900 hectáreas---, el 75, y el que se arma con las parcelas del catastro rural que acompaña las capas, el 74. Sobre el contorno censal, el de escala 1:250.000 y el catastral, entre el 98 y el 99 por ciento del Área de Producción y Conservación queda dentro de la cuenca. \textit{Sobre la capa vigente del IGN ese último cruce no se hizo: se la incorporó después, cuando el ráster del mapa final ya no estaba entre los materiales procesados, y sólo se midió la proporción de cuenca.}

\textbf{Entre una quinta y una cuarta parte del departamento no tiene categoría.} El mapa deja sin categoría entre 22.400 y 26.800 hectáreas, según el contorno ---el sector noroeste y el que corresponde a la cuenca Grande--Perico, que no figura en la Tabla 17--- porque caen fuera de la capa de cuencas con que se construyó. Con la cartografía censal el sector tiene nombre: 17.700 de esas hectáreas son el \textbf{radio rural 0101 de Vaqueros}, donde el censo contó setenta y ocho habitantes (capítulo \ref{cap:poblacion}), y el resto se reparte entre los radios rurales 0201 y 0202 de la fracción de La Caldera. \textbf{Y con la capa catastral el sector tiene dueño}: unas 19.700 de esas hectáreas son el catastro 102, «Las Nieves», \textbf{tierra fiscal de la Provincia} y territorio que reclama la comunidad kolla Kondorwaira (capítulo \ref{cap:resistencias}, lámina \ref{fig:kondorwaira}). El nuevo artículo 14 de la Ley 7543 pone en Categoría II los inmuebles fiscales con cobertura boscosa; el mapa deja sin ninguna categoría el inmueble fiscal más grande del departamento. Si hay bosque nativo en él y qué régimen lo alcanza, las capas no lo dicen; el apéndice \ref{cap:pedidos} lo pregunta. \textit{Las tres superficies por categoría de este apartado ---28.600, 45.500 y 3.100 hectáreas--- suman 77.200 y, sobre el contorno censal de 107.600, dejarían 30.400 sin clasificar: la diferencia con las hectáreas sin categoría, unas cuatro mil, no se explica con las capas que este libro procesó.}
\end{cautela}

El capítulo \ref{cap:conclusion} vuelve sobre la consecuencia práctica de este cruce: que el desmonte de bosque nativo está vedado en casi todo el departamento.

Dicho lo cual, el efecto que este capítulo describe subsiste, y es lo que hay que sostener con precisión.

La opacidad, aquí, no es un propósito: es una \textbf{consecuencia} del diseño. Un instrumento técnicamente sofisticado, participativo y bien fundado produce, para el ciudadano común, exactamente el mismo resultado práctico que produciría uno opaco: no se puede saber si un bosque determinado es desmontable hasta que alguien presente un proyecto sobre él. El dato no está oculto; \textbf{no existe todavía}.

Y hay un desplazamiento del sujeto que sí conviene registrar. La Resolución 00509 establece que será el titular del inmueble quien identifique dentro del APC la Categoría III sobre la cual se establecerá el proyecto. En el mapa clásico, el Estado decide dónde y el privado consulta. En este esquema, el privado propone dónde y el Estado verifica el cupo. La diferencia no es de legalidad ni de honestidad: es de iniciativa.

De modo que la formulación correcta no es que alguien escondió el mapa. Es que la provincia adoptó un instrumento cuya verificabilidad ciudadana depende enteramente de herramientas informáticas de consulta, y que quien no las domina queda, en los hechos, sin poder responder la pregunta más simple sobre el suelo que tiene enfrente.

\section{Dónde se decide realmente}

La suma de las tres capas produce un desplazamiento que este libro encontró por eliminación.

Los actos políticos visibles que el capítulo \ref{cap:politica} reúne ---el voto del senador contra el ordenamiento de bosques, el Concejo que pide clausurar el ingreso a un loteo y abre el perilago, el municipio que denuncia un loteo ilegal, el intendente que reconoce su insolvencia, el jefe de Gabinete que pide más control--- no son los de un poder capturado.

La decisión que produce el resultado no pasa por los órganos políticos visibles. Pasa por el certificado de aptitud ambiental municipal, por el dictamen del Órgano Técnico de Aplicación y por el expediente de la Secretaría de Ambiente provincial. Tres instancias técnicas, sin votación nominal publicada, sin registro accesible en línea, sin cobertura periodística sistemática y sin obligación de publicar sus fundamentos.

Ése es el mismo hallazgo que \textit{El dispositivo salteño} formuló para la escala provincial: el desplazamiento del órgano político visible al órgano técnico opaco. Que reaparezca en escala municipal, con actores distintos de los que aquel libro estudiaba, es el argumento más fuerte a favor de que se trata de un dispositivo y no de una suma de conductas.

\section{Ocho lugares donde nadie mira}

El apartado anterior estableció dónde se decide. Falta el mecanismo: \textbf{por qué el
resultado se produce sin que nadie tenga que decidirlo}. La respuesta que este libro puede
ofrecer es que el circuito tiene lugares donde la decisión no se toma, no se ve o no se
registra, y que esos lugares no están repartidos al azar.

\begin{cautela}
Lo que sigue es \textbf{inferencia, no documento}. Cada uno de los ocho lugares se apoya en
hechos que este libro documentó y que se citan en su fila; el que los reúne y los presenta como
un conjunto con propiedades comunes es este libro, y ningún expediente lo dice así. El último
apartado enuncia qué hallazgo rompería la construcción, que es la única forma de ofrecerla sin
pedir que se la crea.
\end{cautela}

\begin{ficha}{Ocho lugares donde nadie mira}
\begin{enumerate}[leftmargin=*,nosep]
\item \textbf{El mostrador.} El criterio se aplica sin registrarse. La Zona PU se declara no
urbanizable y el artículo 31 del Código la convierte en trámite de Proyecto Especial; el
proyecto de ordenanza del PDUA deja seis decisiones «a consideración del municipio» sin decir
qué se evalúa, con qué dictamen ni en qué plazo (capítulos \ref{cap:cot} y \ref{cap:pdua}). No
hay expediente del criterio porque el criterio se forma en la atención al público.
\item \textbf{El canal.} La información existe, se publica y no viaja por donde mira el
interesado. La línea de ribera sale dos días en el Boletín y no se asienta en la matrícula
(capítulo \ref{cap:ribera}). El que vende ya lo sabe; el que compra tendría que haber leído el
Boletín de un día determinado de hace once años. \textbf{Y la distinción no es una invención de
este libro}: un auto judicial de 1920, publicado en enero de 1921, ordena publicar un edicto de
deslinde \textbf{treinta días en dos diarios de la capital y «\textit{por una sola vez}» en el
Boletín Oficial} (capítulo \ref{cap:siglo}). El juez separaba, en la misma frase, la publicidad
que da noticia de la que deja constancia.
\item \textbf{El reparto.} Dos organismos intervienen y cada uno supone que el otro miró. La
Resolución 646D/21 registra que Recursos Hídricos intervino a fojas 59 y no dice qué certificó
(capítulo \ref{cap:loteo}). La competencia está dividida; la responsabilidad de mirar, no.
\item \textbf{El espejo.} El que promueve es el que controla: promover, autorizar, evaluar e
inspeccionar la extracción es tarea de una misma oficina (capítulo \ref{cap:amparo}). No hace
falta que nadie se distraiga; no hay quien mire desde afuera.
\item \textbf{El aparato ausente.} La norma presupone un ojo que no existe. Un municipio de
tercera categoría que en 2020 declaró no tener con qué pagar sueldos recibió un Código de
ciento cincuenta y dos artículos diseñado para un cuerpo técnico que no tiene (capítulos
\ref{cap:cot} y \ref{cap:hacienda}). El control está previsto y no está instalado.
\item \textbf{El dato que no se produce.} El parcelario que se ve y no se consulta, el archivo
descargable que se corta en 2015 y el índice que no dice de qué trata nada, todos en este mismo
capítulo. No es información restringida: es información que no se produjo en forma consultable.
\item \textbf{El tiempo.} El acto y su consecuencia quedan separados por más de lo que dura una
gestión ---y a veces el acto llega tarde a su propia fecha: la convocatoria que repone la
elección anulada de 1921 en el distrito de la Caldera se publica \textbf{diecinueve días después
del día que ella misma señalaba para votar} (capítulo \ref{cap:siglo}). La línea de los arroyos El Durazno y Guaranguay se determinó en 2014 y esos arroyos
inundaron el loteo en 2018 (capítulos \ref{cap:ribera} y \ref{cap:amparo}); un trámite de
concesión de agua que lleva nueve años atraviesa dos mandatos completos (capítulo
\ref{cap:redes}).
\item \textbf{La escala.} Se mide el departamento y se decide la parcela. El censo describe
poblaciones, el registro individualiza inmuebles, y ningún organismo tiene por tarea cruzarlos:
los cruces de los capítulos \ref{cap:poblacion} y \ref{cap:ribera} este libro tuvo que hacerlos
a mano.
\end{enumerate}
\end{ficha}

\subsection*{Qué hacen, que es distinto de qué son}

\textbf{Tienen dirección.} Un lugar que ciega a las dos partes por igual se corrige solo,
porque le molesta a alguien con capacidad de reclamar. Éstos ciegan a un lado: al comprador y
no al vendedor, al vecino y no al loteador, al que no tiene abogado y no al que lo tiene. Es lo
que convierte una falla en un sesgo con sentido.

\textbf{No eliminan el costo de averiguar: lo reasignan.} Saber por dónde corre una línea de
ribera cuesta lo mismo en los dos casos. Lo que el lugar ciego decide es quién lo paga.

\textbf{No tienen dueño, y por eso no se reparan.} En cada uno, todos cumplieron: el organismo
publicó, el municipio visó, el registro inscribió lo que le presentaron. No hay funcionario
responsable del hueco porque el hueco está entre dos escritorios. \textbf{Ahí está la inercia}:
el dispositivo no se mantiene, se sostiene solo, y sobrevive a los cambios de gobierno, de
partido y de ley porque ninguno de esos cambios toca el lugar donde está.

\textbf{Cuando se los tapa, se desplazan.} La Ley 8126 obliga desde 2018 a los municipios a
publicar sus ordenanzas en el Boletín y en internet. Vaqueros lo cumple por su cuenta y el
Boletín provincial sigue remitiendo a sus oficinas (capítulo \ref{cap:vaqueros}). El hueco no
se cerró: se mudó de la ordenanza al anexo.

\textbf{Se componen, y ninguno alcanza solo.} Ésa es la razón por la que ningún acto individual
resulta irregular y el resultado se produce igual.

\begin{cautela}
De lo anterior se sigue una consecuencia que corre en los dos sentidos y conviene decirla
entera. Si la suma de estos lugares alcanza para producir el resultado, \textbf{una intención
sería redundante}: no hace falta que nadie lo haya querido. Pero exactamente por eso
\textbf{una intención podría alojarse ahí sin dejar rastro}, porque el dispositivo, por
construcción, no registra lo que pasa en ellos. Este capítulo no puede decidir entre las dos
lecturas, y describir el mecanismo no es un argumento a favor de ninguna.
\end{cautela}

\subsection*{Qué rompería esta construcción}

No se ofrece como interpretación sino como algo que el próximo documento puede quebrar. Tres
pruebas concretas, y ninguna exige acceso extraordinario:

\begin{itemize}[leftmargin=*,nosep]
\item \textbf{El canal se cae} si los informes de dominio muestran que las determinaciones de
línea de ribera posteriores a 2017 sí quedaron asentadas en sus matrículas. Entonces el
problema sería del período y no del circuito.
\item \textbf{El reparto se cae} si la foja 59 del expediente 18-37.348/2019-0 contiene un
certificado con contenido y no una constancia de intervención. Entonces alguien miró, y lo que
hay es un registro incompleto y no un hueco de competencia.
\item \textbf{El conjunto entero se cae} si en las fojas de El Durazno, en los antecedentes
dominiales del corredor o en las respuestas a las cartas de réplica aparece una instrucción o
una gestión dirigida. Entonces no hay ocho lugares donde nadie mira: hay alguien que miró y
decidió, y lo que corresponde explicar es una decisión y no un dispositivo (capítulo
\ref{cap:conclusion}).
\end{itemize}

\section{El parcelario que se ve y no se consulta}

Hay un caso más, y es el único de este capítulo que cualquier lector puede verificar por sí
mismo en cinco minutos, sin pedir nada a nadie.

La provincia mantiene una \textbf{Infraestructura de Datos Espaciales} ---IDESA--- con visor
público en \textit{mapas.salta.gob.ar}. El capítulo \ref{cap:amparo} documenta que el informe de mapas de 2022, incorporado al expediente, declara haberla usado como fuente cartográfica.

Las capas catastrales que el visor atribuye a cada organismo se reproducen en la lámina \ref{lam:idesacat}.

\begin{ficha}{Consulta al visor de IDESA, 15 de septiembre de 2026}
Bajo la categoría \textbf{Infraestructura}, el visor publica la capa \textbf{«Catastro Urbano y
Rural»}, con tres clases activables: \textsc{rural}, \textsc{subrural} y \textsc{urbano}.

\textbf{La capa cubre el departamento La Caldera.} Al acercarse sobre el corredor se dibujan los
polígonos parcelarios: el ejido urbano del pueblo, los fraccionamientos del valle y, al este del
cerro Pucheta, dos conjuntos de parcelas de gran extensión.

Bajo el organismo \textbf{Municipalidad de Salta}, y sólo bajo él, figuran además dos capas
fechadas ---\textbf{«Catastros Sep2025»} y \textbf{«Catastros Jun2026»}---, cada una con
\textbf{«Polígonos Catastro»} y \textbf{«Etiquetas Catastro»}.

\textbf{En el visor no se obtuvieron los atributos de ninguna parcela del departamento}: el
clic devuelve la ubicación ---departamento, municipio, provincia--- y no la matrícula, la
nomenclatura catastral ni la superficie.
\end{ficha}

\begin{cautela}
Hay que decir con precisión qué se observó y qué no.

\textbf{Lo observado} es el comportamiento del visor público en una fecha. No se probó todo el
portal, no se consultaron sus servicios de datos ---que pueden existir y exponer lo que el
visor no muestra--- y no se pidió acceso por ninguna vía formal. Es perfectamente posible que
el dato esté disponible y que este relevamiento no haya sabido llegar a él.

\textbf{Lo que sí queda registrado} es la diferencia entre dos capas del mismo portal: el
catastro de la \textbf{capital} se publica con etiquetas y en dos versiones fechadas; el de los
departamentos se publica \textbf{como dibujo}.

\textbf{Y la consecuencia práctica, que es lo que importa.} Este libro necesitó, capítulo tras
capítulo, saber a qué matrícula corresponde un polígono. Lo pidió a la Dirección General de
Inmuebles ---veintiséis veces, según el apéndice \ref{cap:pedidos}--- porque no encontró
manera de obtenerlo de otro modo. \textbf{El polígono está publicado; lo que lo vuelve útil,
no.}

Es la misma forma que este capítulo viene describiendo, ahora en el instrumento más moderno de
todos: no falta el dato, falta la parte del dato que permitiría usarlo.
\end{cautela}

\lamina{lam-idesa-catastros.jpg}{lam:idesacat}%
{Capas de catastro del visor de IDESA bajo el organismo «Municipalidad de Salta»}%
{\textbf{Panel de capas del visor de la Infraestructura de Datos Espaciales de la provincia},
filtrado por organismo. Bajo \textbf{Municipalidad de Salta} figuran dos conjuntos fechados
---\textbf{«Catastros Sep2025»} y \textbf{«Catastros Jun2026»}---, cada uno con
\textbf{«Polígonos Catastro»} y \textbf{«Etiquetas Catastro»}. La capa provincial «Catastro
Urbano y Rural», que cubre los departamentos, no tiene equivalente etiquetado.}%
{Captura del visor \textit{mapas.salta.gob.ar}, 15 de septiembre de 2026. Se recortó la
porción de mapa para dejar legible el panel; no se alteró ningún tono ni rótulo.}

\subsection*{Lo que la ley manda, y que este capítulo no había leído}

Este capítulo analizó la Ley 8483 por sus efectos cartográficos ---las superficies, el APC, el
área verde--- y no por su articulado. El articulado contiene tres disposiciones que alcanzan
directamente a lo que el resto del libro documenta.

\begin{ficha}{Ley provincial Nº 8483, sancionada el 19 de diciembre de 2024, promulgada por
Decreto 896 del 27 de diciembre --- expediente 91-51.703/24}
\textbf{Artículo 5º}, que sustituye el artículo 20 de la Ley 7543: «\textit{La realización de
\textbf{urbanizaciones}, obras públicas o de infraestructura, prospecciones u obras energéticas
o de vías de transporte \ldots\ en las zonas comprendidas en las Categorías establecidas en el
OTBN y que conlleven un cambio de uso de suelo, serán autorizadas por la autoridad competente
\ldots\ debiendo someterse al procedimiento establecido en el Título III Capítulo VI de la Ley
7070 y \textbf{deberá contar, vista previa mediante, con la no objeción de la Autoridad de
Aplicación} de la presente Ley}».

\textbf{Artículo 3º}: los proyectos de Categoría III «\textit{deberán contar con la
\textbf{consulta libre, previa e informada}, consistente con lo previsto en el artículo 3º
inciso 11) de la Ley 7543 y el \textbf{Convenio OIT Nº 169}, debiendo garantizarse asimismo los
principios de igualdad, publicidad, oralidad, informalidad y \textbf{gratuidad} de la
participación}».

\textbf{Artículo 9º}: la Autoridad de Aplicación, «\textit{en un plazo de \textbf{180 días} a
partir de la promulgación}», deberá «\textit{concluir la \textbf{identificación de los
desmontes ejecutados desde la sanción del primer OTBN e identificar su legalidad}, debiendo
determinar un plan de acción \ldots\ a fin de dar cumplimiento a lo establecido en el artículo
40 de la Ley Nacional 26.331}».

Y el nuevo \textbf{artículo 14}: la Categoría II comprende las áreas «\textit{que poseen una
\textbf{pendiente superior al quince por ciento}}» o suelos con limitaciones severas, e
\textbf{«integran esta categoría los inmuebles fiscales con cobertura boscosa de propiedad de la
provincia»}.
\end{ficha}

\begin{cautela}
\textbf{El artículo 5º es el que más importa, y conviene decir por qué.}

Toda urbanización que implique cambio de uso de suelo en zona categorizada por el OTBN
\textbf{requiere la no objeción de la autoridad forestal}, además del procedimiento ambiental de
la Ley 7070. \textbf{Este libro documenta, en el capítulo \ref{cap:loteo}, urbanizaciones
aprobadas en este corredor, y no tiene la no objeción de ninguna.}

No afirma que falte: la norma es de diciembre de 2024 y varios de esos loteos son anteriores, de
modo que el artículo 7º ---que excluye el efecto retroactivo--- los alcanza. \textbf{Lo que sí
queda establecido es el requisito para lo que venga}, y que el apéndice \ref{cap:pedidos} lo pide desde que
este capítulo lo leyó.

\textbf{El nuevo artículo 14 agrega una segunda pieza.} La Categoría II comprende las áreas de
pendiente superior al quince por ciento ---que es, precisamente, el tipo de ladera sobre la que
la nota de prensa de 2026 sobre las obras de Vialidad Nacional registra emprendimientos privados sin autorización, según el
capítulo \ref{cap:vaqueros}--- \textbf{y los inmuebles fiscales con cobertura boscosa}, que es
lo que el capítulo \ref{cap:tierrafiscal} sigue expediente por expediente.

\textbf{Y el artículo 9º fija un plazo que ya venció.} Promulgada la ley el 27 de diciembre de
2024, los ciento ochenta días se cumplieron a fines de junio de 2025. \textbf{Este libro no sabe
si esa identificación de desmontes y de su legalidad se concluyó, ni si se determinó el plan de
acción.} Es un pedido nuevo y de los que este libro prefiere: un deber con fecha cierta y
vencida.
\end{cautela}

\subsection*{Y el mapa que sí es consultable no está actualizado}

Existe un visor de terceros que publica el ordenamiento de bosques de las provincias en formato
consultable. Sobre Salta, advierte lo siguiente:

\begin{ficha}{Nota del visor, consultada el 16 de septiembre de 2026}
«\textit{El Ordenamiento que se muestra en el mapa \ldots\ corresponde a la Ley \textbf{sin la
actualización del 2025}. Al momento \textbf{no se cuenta con la información en el formato que
permita la actualización} de dicha sección}».
\end{ficha}

\begin{cautela}
\textbf{Es el mismo patrón, en su versión más nítida.} La ley está vigente desde enero de 2025,
su cartografía es el Anexo II, la actualización quedó acreditada por disposición nacional en
octubre de ese año. \textbf{Y quien se dedica a publicar estos mapas dice que no puede
actualizar el suyo porque la información no está en un formato que lo permita.}

\textbf{El informe técnico remite a un repositorio digital con las capas del mapa final y a un Visor desde el que pueden descargarse.} Este libro no sabe si esas capas coinciden con el Anexo II aprobado por la ley. El repositorio sí estaba accesible: de él salen las capas de la lámina \ref{fig:otbn}, descargadas en septiembre de 2026. Y no atribuye a nadie la intención de impedirlo: puede ser una limitación
técnica del visor, o un anexo en PDF, o un trámite que nadie hizo.

\textbf{Lo que registra es el efecto, que es el que este capítulo describe en otros seis
lugares:} la norma se publica, el polígono se aprueba, y lo que permitiría saber qué le toca a
cada parcela no circula en forma utilizable. \textbf{Un tercero que quiere publicarlo no puede.}
\end{cautela}

\subsection*{Aun así, el mapa viejo muestra algo}

Consultado sobre el corredor, ese visor devuelve una imagen legible, con la advertencia puesta
de que corresponde al ordenamiento anterior a 2025.

\begin{ficha}{Categorías de conservación sobre el corredor, según la cartografía previa a 2025}
\textbf{El amarillo ---Categoría II, mediano valor--- cubre prácticamente todo el
departamento}, desde el filo occidental hasta el límite con General Güemes.

\textbf{El rojo ---Categoría I, muy alto valor de conservación--- forma bandas angostas que
siguen los cursos de agua}: una baja por el río La Caldera desde el norte hasta el pueblo y
continúa hacia Vaqueros; otra recorre hacia el este, hacia Campo Santo.

\textbf{El verde ---Categoría III, bajo valor--- aparece en manchones}, concentrados al este del
departamento y en algunos sectores próximos al área urbanizada.
\end{ficha}

\begin{cautela}
\textbf{La coincidencia entre el rojo y los cauces es lo que este libro debe registrar, y sólo
eso.}

\textbf{La categoría de mayor protección del ordenamiento de bosques sigue, sobre este
corredor, las mismas márgenes} donde el capítulo \ref{cap:resistencias} documenta veinte
concesiones de áridos y el capítulo \ref{cap:ribera} reconstruye más de dos décadas de
determinaciones de línea de ribera.

\textbf{No hay contradicción normativa en eso}: la Categoría I protege el bosque, no prohíbe la
extracción de áridos, que se rige por el Código de Minería y tramita ante otra autoridad.
\textbf{Son dos regímenes que recaen sobre el mismo suelo y no se consultan entre sí}, que es lo
que el apéndice \ref{cap:propuestas} propone corregir con un padrón único.

\textbf{Y hay que decir tres veces qué NO prueba esta imagen.} Es la cartografía
\textbf{anterior} a la revisión de 2025, según la advertencia del propio visor. Es una captura
de pantalla \textbf{sin escala gráfica}. Y este libro \textbf{no midió superficies ni verificó
límites}: describe lo que se ve. La categoría de cada sector según el informe técnico ya se calculó con esas capas (lámina \ref{fig:otbn}); lo que sigue pedido en el apéndice \ref{cap:pedidos} es la versión ajustada por la adenda y aprobada como Anexo II.
\end{cautela}

\subsection*{Y dos instrumentos del informe que el libro no había registrado}

El resumen público del informe técnico ---elaborado por el equipo que convocó el Decreto
3749/2014 e integrado por \textbf{Lucas Seghezzo, J.\ N.\ Volante, M.\ A.\ Iribarnegaray, J.\
L.\ Agüero, C.\ Ortega Insaurralde, C.\ D.\ Venencia y A.\ L.\ Zelarayán}--- describe dos
mecanismos que este capítulo no había consignado y que valen para el conjunto del libro.

\begin{ficha}{Umbral de deforestación por cuenca}
«\textit{Las limitaciones a la aprobación de áreas verdes por cuenca se vinculan a la historia
productiva de la cuenca. \textbf{A mayor deforestación pasada, mayores restricciones
futuras}}».

«\textit{\textbf{No se autorizarán proyectos de cambio de uso de suelo en cuencas excedidas en
su umbral máximo de deforestación}, calculado mediante el método del Número de Curva}».

Y sobre el efecto buscado: visibilizar relaciones causales «\textit{podría promover acciones
colectivas a nivel de cuencas para optimizar el uso compartido de la cuenca, por ejemplo,
\textbf{la reactivación de los consorcios de cuencas}}».
\end{ficha}

\textbf{Eso es, en materia de bosques, exactamente lo que este libro propone para el agua.} El
apéndice \ref{cap:propuestas} argumenta que la cuenca del río La Caldera necesita un órgano que
la administre como unidad, y señala que el Código de Aguas tiene consorcios de usuarios pero
ninguno de segundo grado constituido. \textbf{El informe del OTBN llega a la misma conclusión
por otro camino y sobre otro recurso}: la cuenca como unidad de gestión, y los consorcios como
la herramienta a reactivar.

\begin{ficha}{Criterio 11 y protocolo de consulta}
El criterio 11 de la Ley 7.543 ---«\textit{valor y uso dado por comunidades indígenas y
campesinas a áreas boscosas o colindantes}»--- «\textit{será una \textbf{capa de control de
gestión}}», contenida en el \textbf{Mapa 11 del Anexo IX}.

«\textit{Cada proyecto de manejo o cambio de uso del suelo \ldots\ será sometido a un
\textbf{protocolo específico de consulta previa, libre e informada} \ldots\ El proyecto
\textbf{no podrá ser aprobado si no obtiene el visto bueno de las Comisiones de Evaluación}
conformadas para la aplicación del protocolo}» ---Anexo VIII---.
\end{ficha}

\begin{cautela}
\textbf{Esto le da al capítulo \ref{cap:resistencias} un instrumento que no tenía.} Ese capítulo
documenta el relevamiento territorial de la comunidad bajo la Resolución INAI 617/2013 y no
sabe qué efecto jurídico tiene hoy. \textbf{El OTBN revisado le da uno concreto}: si el
territorio relevado está en el Mapa 11, todo proyecto de cambio de uso sobre él requiere
consulta previa y visto bueno de una Comisión de Evaluación.

\textbf{Y lo que el mapa muestra} es una mancha rosa ---«potencial para el uso de comunidades
indígenas»--- que cubre el Valle de Lerma desde el rótulo de La Caldera hasta Coronel Moldes
(Ley 8483, Anexo I, Anexo IX, Mapa 11, pág.\ 87). A la escala del escaneo no se puede saber si
el territorio relevado bajo la Resolución INAI 617/2013 está dentro ni qué superficie del
departamento abarca: el mapa se publica como imagen, sin capa descargable. Sigue pedida la capa
digital.

\textbf{Una precisión sobre esta fuente.} Las citas anteriores salen del \textbf{resumen} del
informe técnico publicado en el repositorio institucional del organismo de ciencia; la lectura
del Mapa 11, del Anexo I completo publicado con la ley. Es una de las pocas \textbf{fuentes académicas con autores identificados} que este libro cita; el apéndice \ref{cap:fuentes} las enumera.
\end{cautela}

\section{La omisión tiene forma}

Los apartados anteriores presentaron casos. Éste los mira juntos y trata a la omisión como lo
que este libro sostiene que es: no un accidente que se repite, sino un objeto con propiedades.
A un objeto así se le pueden hacer seis preguntas ---qué, quién, cómo, de qué forma, cuánto y
cuándo---. Este capítulo contesta las cuatro primeras con documentos y declara por qué no
contesta las dos últimas.

\subsection*{Qué se omite}

Siempre lo mismo, y cabe en una línea: \textbf{se publica la norma de atribución y se reserva
la norma de aplicación}. Lo que falta es o bien el parámetro con el que se decide ---la
superficie mínima, la obra exigible, el puntaje de una discapacidad frente a una familia
numerosa, el plazo para impugnar un padrón, las coordenadas de una línea de ribera--- o bien la
identidad de quien recibe ---la nómina de adjudicatarios, el contrato con sus exenciones---.

La auditoría de pendientes del final de este capítulo llega a lo mismo por otro camino: lo
normativo se consiguió casi todo y lo dominial y societario casi nada. Son dos formas de la
misma frase. \textbf{Se publica la regla y se guardan el número y el nombre.}

\subsection*{Quién omite}

No hay un emisor. En el inventario de este capítulo aparecen los dos municipios, la Secretaría
de Recursos Hídricos, el Poder Ejecutivo provincial y el organismo que administra la tierra
fiscal: dos niveles de gobierno y cinco reparticiones. Eso es, justamente, lo que impide leer
el patrón como una decisión de alguien.

Y hay un segundo actor que conviene separar del primero, porque hasta aquí este libro los trató
juntos y no son lo mismo: \textbf{el que publica no es el que dicta}. La edición web que no
ofrece el enlace «Ver anexo» en ninguno de los instrumentos que quedaron en la lista, el
archivo descargable que se corta en 2015 y el índice que no dice de qué trata nada son
omisiones del publicador, no de los organismos que dictaron los actos. Un acto puede haber
acompañado su anexo y el hueco aparecer igual, aguas abajo. \textbf{Distinguir las dos columnas
---quién dicta y quién publica--- es condición para poder medir cualquiera de las dos.}

\textbf{Y hay una tercera columna, que este libro encontró en un solo caso y que conviene dejar
anotada porque no se parece a las otras dos.} El anexo que la Secretaría de Minería y Energía
remitió al juzgado en junio de 2022 ---los tres cuadros con las veinte concesiones de áridos del
departamento--- está firmado por tres funcionarios de la \textbf{Secretaría de Recursos
Hídricos}, y la nota que lo encabeza \textbf{no lleva firma de nadie} (capítulo
\ref{cap:resistencias}; láminas \ref{lam:notamineria} y \ref{lam:iia}). El que dicta, el que
publica y \textbf{el que certifica} pueden ser tres distintos. Cuando lo son, la pregunta «quién
responde por este dato» no tiene respuesta dentro de la propia pieza. Es una forma de la omisión
que no consiste en no publicar: consiste en publicar sin sujeto.

\subsection*{Cómo se omite: seis mecanismos}

\begin{ficha}{Los seis mecanismos, del más neto al más difuso}
\begin{enumerate}[leftmargin=*,nosep]
\item \textbf{El anexo no acompaña al acto.} La Ordenanza 437 de Vaqueros publica cuatro
artículos aprobatorios y no el Reglamento de Fraccionamiento; los Decretos 3069/14 y 463/15
publican la ratificación y no el contrato; las Resoluciones 142/14 y 236/14 aprueban y
rectifican una determinación de línea de ribera sin publicar las coordenadas, ni las erróneas
ni las corregidas.
\item \textbf{El anexo acompaña, pero no es legible ni buscable.} El de la Resolución 82/16
---la cabeza de la cadena del régimen de tierra fiscal, a la que las demás remiten--- está
publicado como copia certificada de seis folios escaneada sin capa de texto: no aparece en
ninguna búsqueda y hay que abrirlo para saber que está.
\item \textbf{El anexo está publicado en otro lugar, que el acto no indica.} El sexto
instrumento de la lista resultó estar en el digesto del municipio que lo dictó, no donde lo
publicó el acto.
\item \textbf{No hay texto que omitir, porque la decisión se toma caso por caso.} La Zona PU se
declara no urbanizable y el artículo 31 del Código convierte esa prohibición en un trámite de
Proyecto Especial; el proyecto de ordenanza del PDUA deja seis decisiones «a consideración del
municipio» sin decir qué se evalúa, con qué dictamen, en qué plazo ni con qué recurso (capítulo
\ref{cap:pdua}).
\item \textbf{El contenido existe, está identificado y cuesta.} Arancel, horario de atención y
mesa de entradas, que para la mayoría de los vecinos del departamento equivalen a un archivo
cerrado.
\item \textbf{El cambio no deja acto.} Algo que el Estado tenía deja de estar, y ningún acto
publicado lo dispone. El departamento se queda sin Comisión Municipal entre julio de 1921 y abril
de 1922, y el archivo lo sabe sólo porque el decreto que se la devuelve la da a los
departamentos «\textit{que en la actualidad carecen de ella}»; entre 1944 y 1946 hay un
interventor ejerciendo sin decreto de designación publicado; y de los cincuenta y dos nombres con
que doce nóminas oficiales dividieron el departamento entre 1909 y 1940, \textbf{veinticuatro no
están hoy en ningún nomenclador del Estado}, sin acto que los suprima (capítulo
\ref{cap:siglo}, lámina \ref{fig:parajesnom}, y capítulo \ref{cap:ausencias}, más adelante).
\end{enumerate}
\end{ficha}

Los seis van de «no está» a «está y no se encuentra», «está y no donde el acto dice», «no
existe porque nunca se escribió», «existe y se paga» y «pasó y no consta». \textbf{El cuarto es
ausencia de documento en sentido propio}, y es a la vez el más interesante y el menos demostrable: probar
que un criterio nunca se escribió exige haber leído todo lo que podría contenerlo.

\textbf{El sexto es ausencia de acto, no de criterio, y es el más demostrable de los seis},
porque el archivo registra el antes y el después: la comisión que en 1914 repartía un río y la
que en 1922 «carece de ella», los topónimos de las nóminas de 1918 y el nomenclador de 2026. Lo
que no se ve es el acto del medio. \textbf{Y lo que lo revela es casi siempre un acto posterior
que registra el efecto sin nombrar la causa}, que es la misma propiedad que la ficha siguiente
llama el acto que certifica su propio hueco.

\begin{cautela}
Con dos reservas. Que el Estado haya dejado de registrar un topónimo no prueba que el lugar haya
desaparecido: puede seguir vivo en el uso, y este libro no fue a preguntarlo (capítulo
\ref{cap:ausencias}, más adelante). Y que un acto no se haya encontrado vale sólo dentro de los tramos leídos
edición por edición; fuera de ellos, puede existir y no haberse buscado bien.
\end{cautela}

\subsection*{De qué forma: la omisión la declara el acto que la comete}

Ésta es la propiedad que ordena a los demás mecanismos ---al sexto, por un acto posterior---, salvo al cuarto, y conviene enunciarla sola.

\begin{ficha}{Once formas del acto que certifica su propio hueco, y ocho son de hace un siglo}
\begin{itemize}[leftmargin=*,nosep]
\item La Resolución 9/16 remite tres veces al plazo de exhibición de padrones de la Resolución
82/16, dos de ellas a su artículo 5º. El número del que depende toda la arquitectura de control
ciudadano del régimen de tierra fiscal está nombrado, ubicado y no publicado como texto
(capítulo \ref{cap:tierrafiscal}).
\item El Decreto 463/15 lleva al pie la leyenda «\textit{anexo para su consulta en oficinas de
nuestra repartición}». El acto dice dónde está lo que no trae.
\item El edicto de la Resolución 383/14 aprueba la determinación de la línea de ribera en ambas
márgenes de dos arroyos. Una determinación de línea de ribera \textit{es} un conjunto de
coordenadas: el edicto aprueba el conjunto y no lo transcribe.
\item \textbf{Y una forma que 1920 repite tres veces en un año: la vacante sin dueño.} Los tres
decretos que mueven el único puesto rentado de policía de partido del departamento ---el 677, que
acepta una renuncia y no dice quién reemplaza; el 724, que nombra «\textit{para ocupar el puesto
vacante}» sin mencionar la renuncia que el propio Boletín publicó tres semanas antes; y el 901,
que declara la vacante sin nombrar al que la dejó--- \textbf{publican el hueco y no su causa}
(capítulo \ref{cap:siglo}). Y el decreto 1140 del mismo año nombra encargado del Registro Civil
«\textit{con antigüedad al 5 de Agosto ppdo.\ fecha desde la cual presta servicios en carácter de
interino}»: \textbf{el acto reconoce dos meses y medio de ejercicio que nunca se publicaron}.
\item \textbf{Y la forma más pura de todas, que el barrido de 1922 incorporó al cierre: el acto
que declara rectificar otro y no dice qué rectifica.} Un decreto del Poder Ejecutivo provincial,
de junio de 1922, rectifica el artículo 2º del decreto 510 «\textit{en cuanto nombra miembros de
la Comisión Municipal de La Caldera}» y a continuación \textbf{repite los mismos cinco nombres}.
Las únicas diferencias visibles entre el texto rectificado y el rectificatorio son el orden de la
enumeración y una inicial (capítulo \ref{cap:siglo}). \textbf{El acto publica que corrigió y no
publica la corrección}: el lector de 1922 tenía los dos textos delante y no podía saber cuál era
la diferencia.
\item \textbf{Y una segunda de ese mismo año, sobre la pieza que más importa, que ya no es un
caso sino una serie.} El decreto 172 del
15 de julio de 1922 \textbf{aprueba} el presupuesto de gastos y cálculo de recursos que la
Comisión Municipal sancionó el 28 de mayo, previo dictamen del Fiscal General, \textbf{y no lo
transcribe}; \textbf{el decreto 1.975 del 14 de noviembre de 1924 hace exactamente lo mismo},
por remisión al expediente 5162 F y sin una sola cifra, para un ejercicio que terminaba en seis
semanas. En las veintisiete ediciones leídas de 1922 y en las cincuenta y dos de 1924 no hay una
sola cifra de la hacienda del departamento. Catorce y dieciséis años antes, el Boletín había
publicado el cuadro completo de entradas y
salidas del mismo municipio, hasta las sillas de la comisión (capítulo \ref{cap:siglo}).
\textbf{Lo que cambió no fue el acceso: fue qué parte del acto se publica}, y para cuando se
publican los dos decretos de aprobación el cambio ya está hecho.

\textbf{Y la serie tiene un final que conviene decir, porque este capítulo casi nunca registra
desmentidos.} En 1946 el Boletín publica el \textbf{balance del movimiento de la renta} de la
comuna del departamento, cuarto trimestre de 1945, \textbf{con sus renglones y sus saldos
mensuales} ---noventa y ocho pesos con quince centavos al cerrar octubre--- y firma al pie de
cada mes. \textbf{Después de siete actos económicos publicados sin una sola cifra entre 1922 y
1944, se publica el número.} Lo firma un interventor, y es el único balance con cifras de todo
el período. \textbf{La opacidad de esta serie no fue un estado permanente: fue una práctica que
empezó, duró veinticuatro años y se interrumpió}, y eso dice más sobre ella que si hubiera sido
constante.
\item \textbf{Y una cuarta forma, que es la más incómoda porque el hueco lo declara el propio
Estado sobre sus propios registros.} El 1º de julio de 1924 el Ministro de Gobierno resuelve
autorizar al Jefe del Registro Civil a «\textit{investigar el destino que hayan tenido los
libros}» de esa dependencia en el departamento de La Caldera, por tratarse de una medida
«\textit{de naturaleza urgente, dado el hecho que la motiva e intereses que el mismo puede
comprometer}», con cooperación de la comisaría (capítulo \ref{cap:siglo}). \textbf{No dice cuál
es el hecho que la motiva, no dice qué intereses, y no publica el resultado}: ningún acto
posterior del año dice qué encontró la diligencia. \textbf{Es el único caso del archivo temprano
en que lo que puede faltar no es un anexo ni una coordenada sino el registro mismo de los
nacimientos y las muertes de un departamento}, y lo que se publica de eso es que hubo que ir a
buscarlo.
\item \textbf{Y una tercera de la misma especie, en 1923, sobre el acto de mayor alcance del
archivo temprano.} El decreto 1161 del 24 de septiembre declara intervenida la Municipalidad de
La Caldera «\textit{vistas las presentaciones de que instruye el expediente N.\ 4541 F}» y
porque «\textit{los hechos anotados}» en ellas hacen procedente la intervención. \textbf{Ni el
expediente, ni las presentaciones, ni cuáles son esos hechos, ni las instrucciones que el
artículo 2º manda dar al Interventor, ni un solo acto del Interventor se publican} (capítulo
\ref{cap:siglo}). \textbf{Lo publicado permite saber que a este departamento le suspendieron el
gobierno y no por qué.}
\item \textbf{Y una quinta, de 1944, que es la forma más pura de todas: el archivo registra el
efecto y no el acto.} El decreto que aprueba el presupuesto municipal de ese año dice
«\textit{visto este expediente en el que el señor \textbf{Interventor de la comuna de la Caldera}
eleva}» el proyecto de ordenanza. \textbf{Hay un interventor ejerciendo, y ningún acto de los
treinta y dos años leídos dice cuándo se lo designó ni quién es}; la intervención sólo termina en
agosto de 1946, cuando un decreto la da por concluida. \textbf{Se puede saber que el gobierno de
este departamento estuvo intervenido al menos tres años, y no se puede saber desde cuándo}
(capítulo \ref{cap:siglo}).
\item \textbf{Y una sexta, que se ve sólo por contraste, y por eso importa.} En 1944 la refección
del local de la comisaría del pueblo dejó \textbf{licitación, adjudicación, contratista con
nombre, certificado final por \$5.018,02, un adicional y la devolución del depósito en garantía}.
En 1945 las obras de defensa del río de la Caldera se llamaron a licitación pública por
\textbf{\$2.515} y terminaron ejecutándose \textbf{«por vía administrativa»}: \textbf{ni
contratista, ni certificado, ni recepción}. \textbf{Del edificio queda el expediente entero; del
río, sólo la autorización.} No es que falte un documento: es que el mismo Estado, el mismo año y
en el mismo departamento, documenta una obra y no documenta la otra.
\item Y uno de hace un siglo, que el barrido de 1918 incorporó al cierre. Los decretos 113 y
190 de ese año nombran los \textbf{Jueces de Agua} de los distritos de Caldera y Calderilla
«\textit{vista la nota que antecede del Sr.\ Comisionado Municipal de Caldera, proponiendo
\ldots\ las personas}». \textbf{La nota que decide quién administra el agua de los dos
distritos está nombrada, ubicada y no publicada}, y los dos decretos nombran personas distintas
para los mismos dos distritos con seis semanas de diferencia, sin que el segundo cite al primero
(capítulo \ref{cap:siglo}).
\end{itemize}
\end{ficha}

De ahí se siguen dos cosas. La primera es \textbf{lo que vuelve legal a la omisión}: el acto
cumple la publicidad que la ley le exige, y lo que queda afuera queda afuera declaradamente. La
segunda es la que sirve para trabajar: \textbf{el universo de omisiones es cerrado y
enumerable, no estimado}. No hay que inferir el agujero. El agujero trae sus propias
coordenadas, y localizarlo es una operación mecánica sobre el corpus ---buscar los actos que
remiten a contenido que no acompañan---, no un juicio sobre lo que debería estar.

El cuarto mecanismo es la excepción y por eso se lo marca aparte: ahí no hay acto que declare
nada, y todo lo que se diga de él es inferencia.

\subsection*{Las dos columnas que este capítulo no escribe}

\textbf{Cuánto.} El único recuento que este libro tiene es el de su propia lista: tenía trece
instrumentos candidatos y \textbf{ocho salieron al verificarse uno por uno}. De los seis que
quedaron, cinco no se consiguieron. Es decir que \textbf{seis de cada diez de lo que parecía
omisión era falla de recuperación} ---del buscador, del repositorio o de quien buscaba---. Ese
número tiene que gobernar cualquier cuantificación futura, y es la razón por la que esta
columna no se escribe todavía: una tasa de omisión calculada sin descontarlo mide sobre todo la
destreza del que busca.

\textbf{Y dos datos más, reunidos al cierre, empujan en la misma dirección.} El capítulo
\ref{cap:ausencias}, más adelante, daba por ausente del archivo la vida asociativa del departamento, y estaba en
el mismo Boletín, bajo un tipo de aviso que el relevamiento no había mirado. Y este libro afirmó
en tres lugares que la directora jubilada en 1918 era el primer nombre de la escuela del pueblo,
cuando otro de sus capítulos transcribía a la de 1911. \textbf{Ninguno de los dos es un anexo,
pero los dos son ausencias que el libro atribuyó al objeto y eran del que buscaba.}

\textbf{Cuándo.} Los instrumentos del inventario se concentran entre 2014 y 2016, que es
exactamente la ventana que este relevamiento leyó con más intensidad y donde además se
concentran los veinticinco actos de ribera del capítulo \ref{cap:ribera}. \textbf{Mientras el
corpus moderno no esté leído entero, la forma temporal de la omisión es indistinguible de la
forma temporal de la lectura.}

\begin{cautela}
Las dos columnas esperan al corpus completo, y entonces hay tres inflexiones concretas que se
podrán probar y que hoy no. La primera es la línea de base que este mismo capítulo estableció:
los anexos sí se publican, y existe una colección digital separada donde el departamento
aparece con código, importe y expediente. La segunda es la Ley 8126 de 2018, que obliga a los
municipios a publicar sus ordenanzas en el Boletín y en internet: hay que ver si bajó la tasa o
si sólo cambió de lugar. La tercera es el corte del archivo descargable en 2015, que es un
cambio del publicador y no de los organismos, y que por eso debería verse en todos los emisores
a la vez.
\end{cautela}

\subsection*{Lo único que hoy puede decirse del tiempo}

En octubre de 1914 el municipio de La Caldera dictó una ordenanza que repartía la totalidad del
caudal del río Santa Rufina y reconocía derechos adquiridos por uso tradicional. Este libro
leyó las setenta ediciones del Boletín de ese año, una por una, y \textbf{el texto no
está}. En enero de 2015, el municipio de Vaqueros publicó una ordenanza que aprueba en cuatro
artículos un Reglamento de Fraccionamiento de Tierras, y \textbf{el reglamento tampoco está}.

Cien años y un departamento de por medio, la misma figura: \textbf{se conoce que la regla
existe porque otro documento la nombra, y no se conoce la regla}.

\begin{cautela}
Los dos casos no tienen el mismo valor probatorio y no deben leerse como una serie. El de 2015
es una omisión documentada: el acto está publicado, se ve qué publica y se ve qué no. El de
1914 es una ausencia, y una ausencia admite dos explicaciones que este libro no puede separar
---que la ordenanza no se haya publicado en el Boletín provincial, cosa esperable en un acto
municipal de la época; que se haya publicado y no se conserve; o que esté donde este libro
todavía no fue a buscarla, que es lo más probable: las fojas del expediente de 1938 donde el
municipio la invocó, o el libro de actas de la comisión municipal (capítulo
\ref{cap:expropiacion})---. \textbf{Son dos puntos, no una tendencia.} La tendencia, si existe, la dirán las dos columnas que este capítulo dejó sin
escribir.
\end{cautela}

\subsection*{Y dos puntos más, que corren en contra}

Los barridos de 1918 y 1919, incorporados al cierre, agregan algo que este capítulo debe poner
aunque no lo favorezca.

\textbf{El decreto 1.660 del 24 de abril de 1918 publica su propia cuenta.} Es el acto que
devuelve a La Caldera la condición de Comisión Municipal, después de que el decreto 1.557 la
hubiera convertido en Municipalidad electiva siete semanas antes, y para fundarlo el Boletín
imprime las tres cifras: el cálculo de recursos de \textbf{\$6.471,44}, el saldo del ejercicio
anterior indebidamente computado, \textbf{\$3.083,94}, y la renta real resultante,
\textbf{\$3.387,50}, contra el mínimo de cinco mil que la norma exigía (capítulo
\ref{cap:siglo}).

\begin{cautela}
\textbf{Es exactamente lo contrario de lo que este capítulo describe para 2014--2026.} Ahí el
acto publica la atribución y reserva el parámetro; acá el acto publica el parámetro, el umbral,
la cuenta y el resultado, en el mismo renglón en que decide. Un vecino de 1918 que abriera el
Boletín podía saber por qué su departamento no se gobernaba a sí mismo, y por cuánto. Un vecino
de 2015 que abriera el Boletín no podía saber qué obras le eran exigibles al loteador de
enfrente.

\textbf{Y no es un caso aislado: el año siguiente lo repite otro gobierno.} El decreto 33 del 15
de enero de 1919, que nombra la Comisión Municipal del departamento, publica el criterio entero
en su considerando ---qué artículo de la Ley Orgánica se invoca, qué artículo de la Constitución
lo derogó, y cuál es el umbral: «\textit{un centro urbano de cinco mil habitantes, de que carecen
todos los Departamentos de Campaña}»--- y hasta declara el programa de gobierno que lo
enmarca (capítulo \ref{cap:siglo}). \textbf{Dos actos, doce meses, dos gobiernos distintos, y
los dos publican la regla con la que deciden}: uno el número, el otro la norma.

\textbf{Dos casos tampoco son una serie}, y conviene decir con qué se los puede tumbar: bastaría
con un acto de la misma época que reserve su criterio para que el contraste se disuelva. \textbf{Y
este mismo capítulo tiene uno}: el decreto 1161 de 1923, que interviene el municipio por «\textit{los
hechos anotados}» en un expediente que no se publica, es de cinco años después y decide sobre lo
mismo ---quién gobierna el departamento---. De modo que el contraste no separa dos épocas: \textbf{en
1918 y 1919 el Estado publicó la regla con la que decidía la condición institucional del
departamento, y en 1923 decidió lo mismo sin publicarla}. Este libro, además, no leyó de corrido
el segundo semestre de 1922 ni nada entre 1947 y 2012 ---1947 a 1950 se barrieron sin lectura
sobre la imagen---, de modo que \textbf{la afirmación fuerte ---que la opacidad de criterio es algo que
apareció y no una constante del Estado provincial--- sigue sin poder hacerse}. Lo que sí queda
establecido es que no siempre fue así, y que hay dos actos fechados que lo prueban.
\end{cautela}

\subsection*{Y una omisión de 1919 que no es de criterio sino de sujeto}

Hay en 1919 una forma que este capítulo no tenía y que conviene anotar, porque no consiste en
reservar un contenido sino en publicar mal quién decide.

\textbf{La mancheta de las cuarenta y nueve ediciones del año atribuye el gobierno al Dr.\
Joaquín Castellanos.} Y sin embargo el Ejecutivo lo ejerce el \textbf{Presidente del Senado} en
treinta y dos de ellas, en dos interinatos separados por un regreso del titular, \textbf{sin que
se publique un solo decreto de apertura ni de cierre de ninguno de los dos} (capítulo
\ref{cap:siglo}).

\begin{cautela}
\textbf{El envoltorio dice una cosa y la firma dice otra, en la misma hoja.} Quien lea la tapa
sabe quién gobierna; quien lea el pie del decreto sabe quién firmó; y nadie puede saber desde
cuándo ni hasta cuándo, porque el acto que lo dispone no se publica.

Es la misma estructura que este capítulo describe, corrida un lugar: no falta el criterio, falta
\textbf{la constancia de quién estaba habilitado para aplicarlo}. Y se parece a lo que el
apartado anterior llama la tercera columna: el que dicta, el que publica y el que certifica
pueden no coincidir, y aquí no coinciden ni siquiera el que gobierna en la tapa y el que firma
adentro.
\end{cautela}

\section{El límite de lo que se puede afirmar}

Corresponde ahora la parte más incómoda, y va sin atenuantes.

Todo lo que este libro documentó tiene una explicación estructural completa y suficiente. Un municipio clasificado en tercera categoría, que en 2020 declaró no tener con qué pagar sueldos, recibió en 2015 un plan urbano y en 2016 un Código de ciento cincuenta y dos artículos diseñado para un aparato técnico que no posee, y absorbió entre 2010 y 2022 el cuarto crecimiento poblacional de la Argentina. La Provincia y el Consejo Federal de Inversiones financiaron la redacción de la norma; nadie financió la estructura que debía aplicarla. El mercado ocupó el espacio que esa incapacidad dejó libre.

Esa explicación no requiere que nadie haya decidido producir el resultado.

\begin{cautela}
Pero la explicación estructural es también, exactamente, lo que un dispositivo produciría si funcionara bien.

Este libro no puede distinguir, con los documentos que consiguió, entre un Estado desbordado y un Estado desbordado que a alguien le conviene que siga desbordado. Los dos se ven igual desde afuera. Un municipio sin capacidad técnica y un municipio al que no se le da capacidad técnica producen el mismo expediente vacío.

Distinguirlos requeriría lo que no está en ningún archivo: saber qué se conversa en las oficinas donde la excepción se tramita. Ese es el límite de una investigación documental, y el libro llega hasta ahí y se detiene.

Lo que en efecto puede afirmarse sin ambigüedad es que la opacidad de este departamento no es un secreto guardado. Es una capa fina y enteramente legal, hecha de anexos no publicados, de excepciones que se tramitan y de actos individuales que nunca se registran públicamente, y basta para que nadie pueda reconstruir una decisión sin ir físicamente a una oficina. Es más eficaz que cualquier ocultamiento, porque no deja nada que denunciar.
\end{cautela}

\section{Las tres cosas que quedaron sin resolver}

No son misterios. Son trámites, y su resolución no depende de encontrar nada oculto.

\textbf{Primero.} Si la actual titular del organismo provincial que administra la tierra fiscal es copropietaria de un loteo cerrado en el departamento que representó ocho años en la Legislatura. Se resuelve con el informe de dominio de una matrícula.

\textbf{Segundo.} Cuándo, a instancia de quién y por qué motivo la sociedad titular del loteo El Durazno fue incorporada como tercero al amparo colectivo que litiga por la inundación de ese loteo. Se resuelve con el expediente completo.

\textbf{Tercero.} Si veintiséis de los treinta adjudicatarios del barrio Juan Manuel de Rosas son vecinos de La Caldera, como sostuvo el centro vecinal. Se resuelve con el domicilio de los adjudicatarios, que ninguna de las dos versiones del anexo consigna.

Las tres se responden con documentos que existen y a los que este relevamiento no accedió. Que sigan sin respuesta no es prueba de nada: es la medida exacta de hasta dónde llegó este trabajo.

\section{La forma del hueco}

Al cerrar el libro se auditaron uno por uno los marcadores de pendiente que quedaron en el texto, para establecer cuáles podían resolverse en fuentes abiertas y cuáles no. El resultado de esa auditoría es, en sí mismo, un dato sobre el objeto, y por eso se consigna aquí y no en un apéndice de método.

Lo que se pudo cerrar es casi todo normativo. Se obtuvieron el texto de la ley orgánica de municipalidades, el de la ley de áreas protegidas, la secuencia completa del régimen nacional de expansión de redes de gas, la identificación de la ley general de expropiaciones, el contenido resolutivo de la última audiencia del amparo y la ubicación exacta de la desagregación censal por municipio.

Lo que resistió es casi todo dominial y societario. Ningún antecedente de finca, ningún informe de dominio, ningún legajo de sociedad; de los expedientes judiciales, piezas del incidente y de la ejecución ---providencias, actas e informes--- pero no la sentencia ni el expediente principal; y el único plano de mensura obtenido ---el 000962--- llegó como anexo de un decreto.

\textbf{Esa asimetría no es un accidente del método de búsqueda: es la arquitectura de acceso del sistema.} Toda la información que permite saber cuánta gente vive en el departamento y en qué condiciones es gratuita, remota y descargable. Toda la información que permite saber quién es dueño de qué es arancelada, presencial y, en algunos casos, requiere acreditar interés legítimo.

\textbf{Y lo dominial no le resiste sólo al vecino: también le resistía al Estado.} En 1939 la Curia promovió posesión treintañal (capítulo \ref{cap:fincas}); en 1943 la Provincia tuvo que promover juicio de posesión treintañal para titular el terreno donde su propia comisaría funcionaba desde 1913 o antes (capítulo \ref{cap:presencia}); y en 2015 adquirió por prescripción el suelo de la escuela de El Gallinato, que funcionaba en terreno ajeno desde 1968 (capítulo \ref{cap:tierrafiscal}). Entre 1939 y 2015, tres ocupantes institucionales tuvieron que titular por prescripción lo que ya usaban. \textbf{El registro de la propiedad no es opaco sólo para quien pregunta desde afuera: tampoco le alcanzaba al Estado para saber qué era suyo.}

\textbf{Y la asimetría tiene una medida exacta en un solo año.} En 1920 el Boletín publica \textbf{diez resúmenes del movimiento de la Tesorería General}, uno por mes y encadenados por la existencia de caja, y \textbf{ningún balance municipal}: ni de este departamento ni de ningún otro (capítulo \ref{cap:siglo}). Diez a cero, en el mismo boletín y en el mismo año. El Estado provincial rinde su caja mes a mes; los municipios, que son los que deciden sobre el suelo, no rinden ninguna, y el capítulo \ref{cap:hacienda} muestra que un siglo después la información fiscal municipal sigue siendo lo más difícil de conseguir de este territorio.

Un Estado que publica generosamente las reglas y resguarda los nombres no está ocultando nada: está distribuyendo el costo de averiguar. Y ese costo lo paga entero quien quiera reconstruir la propiedad del suelo, mientras no lo paga nadie que quiera describir la pobreza.

Es la misma capa fina que este capítulo viene describiendo, vista desde el otro lado. No hay un archivo cerrado. Hay un arancel, un horario de atención y una mesa de entradas, que para la mayoría de los vecinos del departamento equivalen a un archivo cerrado.

\subsection*{Lo que el Estado sí declara de interés}

Una nota breve, y no es ironía.

El 13 de noviembre de 2015, por \textbf{Decreto Nº 3760}, el Poder Ejecutivo declaró de \textbf{Interés Provincial la Fiesta Patronal de Nuestra Señora de la Candelaria} de la localidad de La Caldera, a celebrarse el 2 de febrero de 2016.

Se consigna porque forma parte del registro y porque contradice, en su pequeñez, una lectura fácil de este libro.

El Estado provincial no ignora a La Caldera. La nombra, la visita, declara de interés su fiesta patronal, contrata a su municipio para mantener rutas, le construye aulas, le cede el arancel de las licencias de conducir, sostiene allí centros de mediación comunitaria. La relación es densa y continua.

Lo que este libro describe no es abandono. Es \textbf{selectividad}: un Estado muy presente en algunas dimensiones --- la turística, la vial, la ceremonial --- y sistemáticamente tardío en otras --- el agua de red, la cloaca, el control del fraccionamiento, la publicación de las coordenadas ---. Confundir las dos cosas sería un error de diagnóstico, y de él se sigue un error de reclamo: no se le pide a este Estado que llegue, se le pide que llegue con lo otro.

\subsection*{El río como frontera de la notificación}

Hay una pieza de 1975 que sigue rigiendo y que conviene poner junto a las demás.

\begin{ficha}{Acordada Nº 4529/75 de la Corte de Justicia, citada en el B.O.\ Nº 19.256 de 2014}
La zona en la que la oficina de Ujiería del Poder Judicial realiza las diligencias de notificación está comprendida dentro de estos límites: ``\textbf{Norte: Río Vaqueros}; Sur: Río Ancho; Este: kilómetro 5 de la Ruta Nacional 34; Oeste: río San Lorenzo''.
\end{ficha}

El río Vaqueros vuelve a aparecer como límite, y esta vez no de una competencia municipal sino de la notificación judicial.

Al norte de ese río --- es decir, en Vaqueros y en La Caldera --- las diligencias no las hace la ujiería. Quien litiga desde el departamento debe constituir domicilio \textbf{dentro} del radio, en la capital, y si no lo hace las notificaciones se tienen por practicadas por ministerio de la ley: corren los plazos aunque no se haya enterado.

No es una norma dirigida contra el departamento. Es un radio de ujiería trazado en 1975 sobre accidentes geográficos, cuando al norte del río Vaqueros vivía una fracción de la población actual. Hoy viven doce mil trescientas personas.

Pero el efecto acumulado es el que este capítulo describe. El vecino que quiere apelar una línea de ribera tiene diez días hábiles y las coordenadas en una sede de la capital. El que quiere consultar un plano paga arancel y va presencialmente. El que litiga debe constituir domicilio en la ciudad o asumir que los plazos corren sin él. Ninguna de esas tres reglas fue escrita para excluir a nadie. Las tres tienen el mismo efecto, y las tres se apoyan en la misma premisa tácita: que el interesado está del otro lado del río.

\subsection*{Y a veces no hay nada}

Conviene registrar también el reverso, porque el relevamiento sistemático del Boletín Oficial permite medirlo.

Entre el 11 de junio y el 3 de julio de 2014 --- quince ediciones, quinientas cuarenta y dos páginas --- el departamento La Caldera aparece \textbf{una sola vez}, y es el mismo aviso repetido cinco días por mandato legal: la solicitud de un particular para extraer agua de un pozo. Ni Vaqueros, ni La Calderilla, ni Yacones, ni Lesser figuran en todo el período.

No hay que forzar la lectura. Tres semanas sin actos referidos a un departamento que entonces tenía unos nueve mil habitantes es lo normal en el boletín de una provincia de alrededor de un millón trescientos mil, y la ausencia de noticias no significa ausencia de gestión: significa que los actos que la administración publica --- decretos, licitaciones, edictos --- se producen a un ritmo irregular.

Y el fenómeno se repite y se puede medir mejor. Entre el 10 de noviembre y el 2 de diciembre de 2014 --- dieciséis ediciones, seiscientas veinticuatro páginas --- la expresión ``La Caldera'' \textbf{no aparece una sola vez}. Tampoco ``La Calderilla'', ni ``Yacones'', ni ``río Caldera''. El relevamiento que arroja ese resultado validó la capa de texto de los archivos corriendo reconocimiento óptico propio sobre una muestra y comparando volúmenes de caracteres, y corrió además una búsqueda difusa tolerante a los errores típicos de escaneo: los seis aciertos fueron falsos positivos --- ``escalera comunitaria'', ``zona tabacalera del Río Toro'', ``calle Rafael Obligado''. \textbf{El cero es real, no un artefacto de extracción.} Es, con todo, un cero del término y no del departamento: en esa misma ventana el Boletín Nº 19.430, del 19 de noviembre, publica la convocatoria a asamblea de la asociación vecinal de los ríos Castellanos y Lesser (capítulo \ref{cap:resistencias}).

Lo que puede decirse, y es modesto, es que \textbf{la presencia del departamento en el registro oficial es intermitente y desigual}. Se concentra en rachas --- veinticinco actos hídricos en dos años, seis concesiones de áridos en dos --- y desaparece por semanas enteras, con actos sueltos fuera de las rachas, como la línea de ribera del río Wierna de noviembre de 2016, cuyas coordenadas tampoco se publicaron. Quien quiera reconstruir la vida administrativa del corredor leyendo el Boletín no encontrará una serie continua sino un pulso: períodos densos separados por vacíos.

Eso importa para el método de este libro más que para su tesis. Explica por qué la reconstrucción documental produce imágenes nítidas de algunos años y ninguna de otros, y por qué el silencio de un período no autoriza a concluir que en ese período no pasó nada.

\subsection*{El archivo descargable se corta en 2015}

Este libro se construyó, en buena parte, descargando ediciones completas del Boletín Oficial y leyéndolas. Ese método tiene un límite y conviene dejarlo escrito, porque no es una limitación del trabajo sino del repositorio.

\begin{ficha}{Comprobación del 7 de septiembre de 2026}
El sitio del Boletín Oficial de Salta publica sus ediciones en formato descargable bajo una dirección predecible, del tipo \texttt{/pdfs/<año>/<número de edición>.pdf}.

Probadas cinco ediciones espaciadas de cada año entre 2005 y 2026, el resultado es inequívoco: \textbf{todos los años de 2005 a 2015 responden; todos los de 2016 a 2026 devuelven error 404}. No es un bloqueo por exceso de pedidos --- eso daría otro código --- sino ausencia de los archivos en esa ruta.
\end{ficha}

\begin{cautela}
De modo que la última década del Boletín Oficial \textbf{no se puede descargar en bloque}. Se consulta de a un acto por vez, a través de la interfaz del sitio, y para llegar a un acto hay que saber de antemano su número o su fecha.

Esto no es una denuncia y no debe leerse como tal. Puede obedecer a una migración de plataforma, a una decisión de almacenamiento o a una política de acceso deliberada y perfectamente legítima. Este libro no sabe cuál de las tres, y no lo averiguó.

Lo que corresponde registrar es el \textbf{efecto}, que es idéntico al de los otros casos de este capítulo. La información está publicada: cualquiera puede leer el Boletín de 2019 acto por acto. Lo que no puede es \textbf{leerlo entero}, ni buscar dentro de él, ni preguntarle cuántas veces aparece un topónimo, una empresa o una matrícula. Y esas son, exactamente, las preguntas que permiten reconstruir un dispositivo.

El resultado práctico es una asimetría curiosa: \textbf{los años más antiguos son los más accesibles}. Este libro pudo reconstruir con detalle el ciclo hídrico de 2013 a 2015 --- veinticinco actos, más de veinticinco inmuebles, trece cursos --- y no puede hacer lo mismo con el período que va del primer aluvión a la sentencia. Para 2018 y 2019, los años decisivos de esta historia, la reconstrucción documental depende de conocer de antemano qué buscar.

Quien quiera contar lo que pasó en el corredor entre 2016 y 2026 deberá pedirlo pieza por pieza, y sólo encontrará lo que ya sospechaba.
\end{cautela}

\subsection*{Un índice que no dice de qué trata nada}

Hay un caso que conviene mostrar entero, porque es pequeño, es verificable y explica la arquitectura mejor que cualquier generalización.

El Boletín Oficial de Salta publica en línea sus \textbf{índices de decretos} por año, desde 1983. Se descargan sin arancel, sin trámite y sin identificarse. Es, en apariencia, transparencia retrospectiva completa.

\begin{ficha}{Índice de Decretos, año 1983 --- estructura}
El documento tiene sesenta y siete páginas y tres columnas: \textbf{número de decreto}, \textbf{fecha de publicación} y \textbf{número de edición} del Boletín. Enumera 2.176 decretos, y a partir del 19 de diciembre de 1983 ---así en el índice; el cambio de gobierno fue el 10--- la numeración vuelve a empezar en 1.

No consigna materia, organismo, beneficiario ni objeto. Muchas filas ni siquiera tienen fecha: sólo el número.
\end{ficha}

\begin{cautela}
Léase lo que eso significa para cualquiera que busque algo.

Este libro necesita, para el capítulo \ref{cap:politica}, los decretos por los cuales se designaron los titulares del ejecutivo municipal de La Caldera entre 1983 y 1987 --- presidentes de Comisión Municipal que no se elegían sino que los nombraba el gobernador. El índice permite saber que en 1983 hubo 2.176 decretos y en qué edición del Boletín salió cada uno. \textbf{No permite saber de qué trata ninguno.}

Encontrar uno exige abrir las ediciones del Boletín y leerlas. Para un solo año, eso son más de doscientas cincuenta ediciones. Para el período completo, cuatro veces eso.

El dato no es que la información esté oculta: \textbf{está publicada}. El dato es que está publicada de una manera que sólo sirve a quien ya sabe el número del decreto que busca --- es decir, a quien no necesita el índice. Un índice sin materia no es un índice: es una lista.

Y ésa es, en miniatura, la forma completa del problema que este capítulo describe. La opacidad de este dispositivo casi nunca consiste en negar un documento. Consiste en entregarlo de un modo que preserva su inaccesibilidad práctica: el anexo que no acompaña al acto, las coordenadas que quedan en la sede, el plano que se consulta con arancel, el índice que no dice de qué trata nada.
\end{cautela}
